\documentclass[12pt,a4paper]{article}

% Кодировка и русский язык
\usepackage[utf8]{inputenc}
\usepackage[T2A]{fontenc}
\usepackage[russian]{babel}

% Математические пакеты
\usepackage{amsmath,amssymb,amsfonts}

% Разметка страницы
\usepackage{geometry}
\geometry{
  a4paper,
  left=2cm,
  right=2cm,
  top=2.5cm,
  bottom=2.5cm
}

% Дополнительно: абзацы, межстрочный интервал при желании
\setlength{\parindent}{1.25cm}
\setlength{\parskip}{0pt}

\begin{document}

\begin{center}
    {\Large Решения задач варианта 1}\\[1em]
\end{center}


\section*{Задача 1}

Вы провалились на дно ущелья, глубина от дна до поверхности --- ровно $30$ блоков. Каждый день вы строите столб из булыжника вверх: днём успеваете поставить $6$ блоков, а ночью снимаете $2$ верхних блока и прячетесь внутри, опускаясь на $2$ блока вниз. Утром снова продолжаете с той отметки, где пережидали ночь. Нужно найти, через сколько полных дней столб достигнет поверхности.

За один день происходят два действия:
\begin{itemize}
  \item днём подъём на $6$ блоков;
  \item ночью спуск на $2$ блока.
\end{itemize}
Итого \emph{чистый} подъём за один полный день составляет $6 - 2 = 4$ блока.

Обозначим высоту столба в начале $n$-го дня (до дневного строительства) через $h_n$. Изначально $h_1 = 0$. После первого дня и ночи высота будет
\[
h_2 = h_1 + 4 = 4.
\]
После второго дня и ночи:
\[
h_3 = h_2 + 4 = 8,
\]
и так далее. В общем виде высота после $n$ полных дней (то есть в начале $(n+1)$-го дня) равна
\[
h_{n+1} = 4n.
\]

Однако важно заметить: чтобы выбраться, достаточно, чтобы \emph{днём} какого-то дня, до наступления ночи, столб достиг или превысил отметку $30$ блоков. Ночью вы уже будете на поверхности и не станете прятаться в столбе.

Посмотрим, что происходит в течение $n$-го дня. В начале дня высота равна $h_n = 4(n-1)$. За день вы поднимаетесь ещё на $6$ блоков, то есть дневная высота:
\[
h^{\text{день}}_n = h_n + 6 = 4(n-1) + 6.
\]
Нужно, чтобы
\[
4(n-1) + 6 \ge 30.
\]
Решим это неравенство:
\[
4n - 4 + 6 \ge 30 \quad \Rightarrow \quad 4n + 2 \ge 30 \quad \Rightarrow \quad 4n \ge 28 \quad \Rightarrow \quad n \ge 7.
\]

При $n=7$ дневная высота:
\[
h^{\text{день}}_7 = 4 \cdot 6 + 6 = 30.
\]
Значит, \emph{на седьмой день} вы достигаете поверхности, причём именно днём, до ночного спуска.

\subsection*{Ответ}
В бланк: $7$.

\newpage

\section*{Задача 2}

Рассмотрим величины --- сумму алмазов по чётным и нечётным сундукам отдельно. Обозначим:
\[
S_{\text{чёт}} = \text{сумма алмазов в сундуках с чётными номерами},
\]
\[
S_{\text{нечёт}} = \text{сумма алмазов в сундуках с нечётными номерами}.
\]

Изначально оба эти числа равны $0$. Посмотрим, как одно нажатие на кнопку влияет на $S_{\text{чёт}}$ и $S_{\text{нечёт}}$. Выбираются два соседних сундука. В круге из $16$ сундуков каждый сундук имеет двух соседей: слева и справа. Номер соседнего сундука всегда отличается на $1$, поэтому:
\begin{itemize}
  \item если один сундук имеет нечётный номер, соседний будет чётным;
  \item если один сундук чётный, соседний будет нечётным.
\end{itemize}
Это означает, что любые два соседних сундука состоят из одного чётного и одного нечётного номера.

При нажатии на кнопку:
\begin{itemize}
  \item в один чётный сундук добавляется $1$ алмаз, то есть $S_{\text{чёт}}$ увеличивается на $1$;
  \item в один нечётный сундук добавляется $1$ алмаз, то есть $S_{\text{нечёт}}$ увеличивается на $1$.
\end{itemize}
В результате после каждого нажатия выполняется
\[
S_{\text{чёт}} - S_{\text{нечёт}} \quad \text{не изменяется},
\]
поскольку оба увеличиваются на $1$, и их разность остаётся прежней.

Изначально $S_{\text{чёт}} = 0$ и $S_{\text{нечёт}} = 0$, значит
\[
S_{\text{чёт}} - S_{\text{нечёт}} = 0
\]
на всех шагах работы механизма.

Теперь посмотрим, как выглядит $S_{\text{чёт}}$ и $S_{\text{нечёт}}$ в целевой конфигурации «пароля». В ней в сундуке номер $k$ лежит $k$ алмазов. Чётные номера: $2,4,6,\dots,16$, нечётные: $1,3,5,\dots,15$. Тогда
\[
S_{\text{чёт}} = 2 + 4 + 6 + \dots + 16 = 2(1 + 2 + \dots + 8) = 2 \cdot \frac{8 \cdot 9}{2} = 72,
\]
\[
S_{\text{нечёт}} = 1 + 3 + 5 + \dots + 15 = 8^2 = 64
\]
(последняя формула --- стандартный результат для суммы первых $n$ нечётных чисел).

Разность:
\[
S_{\text{чёт}} - S_{\text{нечёт}} = 72 - 64 = 8.
\]

Но мы уже показали, что в процессе работы механизма $S_{\text{чёт}} - S_{\text{нечёт}}$ \emph{всегда равно нулю} (поскольку стартовое состояние $0$ и разность не меняется). Получить разность $8$ невозможно.

Значит, целевая конфигурация сундуков не может быть достигнута описанным механизмом.

\subsection*{Ответ}
В бланк: Нет.\\
Пояснение: при каждом нажатии на кнопку сумма алмазов в чётных и нечётных сундуках увеличивается одинаково, поэтому их разность всегда равна $0$, а в целевой конфигурации эта разность равна $8$, что невозможно.

\newpage


\section*{Задача 3}

В сундуке лежит $100$ алмазов. Стив и Алекс по очереди забирают из сундука $1$, $2$, $3$ или $4$ алмаза, Стив ходит первым. Выигрывает тот, кто забирает последний алмаз. Оба играют идеально. Нужно определить, кто победит и как именно должен играть победитель.

Это классическая игра на вычитание с конечным числом камней. В таких задачах часто важны остатки по модулю некоторого числа. Заметим, что за один ход игрок может забрать от $1$ до $4$ алмазов, то есть суммарно за два хода (ход Стива и ответный ход Алекса) можно убрать от $2$ до $8$ алмазов.

Главная идея: подобрать стратегию, при которой после пары ходов (ход Стива и ход Алекса) общее количество убранных алмазов всегда будет равно $5$, чтобы после каждого полного цикла ходов (двух ходов) оставалось число, кратное $5$. Для этого можно рассматривать игру по модулю $5$.

Предположим, что Стив стремится после своего хода оставлять в сундуке количество алмазов, кратное $5$. Тогда, если Алекс забирает $a$ алмазов (от $1$ до $4$), Стив на следующем ходу забирает $b$ алмазов так, чтобы
\[
a + b = 5.
\]
То есть:
\begin{itemize}
  \item если Алекс забрал $1$ алмаз, Стив забирает $4$;
  \item если Алекс забрал $2$, Стив забирает $3$;
  \item если Алекс забрал $3$, Стив забирает $2$;
  \item если Алекс забрал $4$, Стив забирает $1$.
\end{itemize}
В результате за каждый «двойной» ход (Алекс + Стив) из сундука уходит ровно $5$ алмазов, и, следовательно, после каждого хода Стива число оставшихся алмазов снова становится кратным $5$.

Изначально $100$ алмазов --- это число, кратное $5$ ($100 = 20 \cdot 5$). Стив делает первый ход, и ему достаточно \emph{в первый ход} снять количество алмазов так, чтобы после его хода осталось число вида $5k + 1$ или $5k$, откуда он сможет придерживаться описанной схемы. Один из удобных вариантов стратегии:

\begin{itemize}
  \item В первый ход Стив забирает $1$ алмаз, остаётся $99$ алмазов.
  \item После этого Стив всегда действует так, чтобы сумма алмазов, забранных за один ход Алекса и следующий ход Стива, была равна $5$.
\end{itemize}

Посмотрим, что будет происходить. После первого хода Стива осталось $99$ алмазов. Далее идут пары ходов (Алекс + Стив), в каждой паре уходит $5$ алмазов. Всего таких пар получится
\[
\frac{99}{5} = 19{,}8,
\]
но корректнее рассуждать по шагам: каждый раз после хода Стива количество алмазов уменьшается на $5$ по сравнению с предыдущим его ходом. Важно, что в конце игры, когда в сундуке останется $5$ алмазов перед ходом Стива, Алекс на своём предыдущем ходу заберёт что-то от $1$ до $4$ алмазов, а Стив заберёт оставшиеся так, чтобы суммарно выйти на $5$ и при этом забрать последний алмаз.

Например, когда осталось $5$ алмазов перед ходом Алекса:
\begin{itemize}
  \item если Алекс забирает $1$ алмаз (остается $4$), Стив забирает $4$ и выигрывает;
  \item если Алекс забирает $2$ (остается $3$), Стив забирает $3$ и выигрывает;
  \item если Алекс забирает $3$ (остается $2$), Стив забирает $2$ и выигрывает;
  \item если Алекс забирает $4$ (остается $1$), Стив забирает $1$ и выигрывает.
\end{itemize}

В любом случае последний алмаз забирает Стив.

Существует и более компактное рассуждение: можно показать, что позиции, в которых число алмазов кратно $5$, являются выигрышными для игрока, который ходит \emph{первым из пары ходов}, и что Стив может всегда возвращать игру в такую позицию после своего хода. Поскольку стартовое число $100$ кратно $5$, Стив при идеальной игре гарантированно победит.

\subsection*{Ответ}
В бланк: Стив.\\
Пояснение: Стив выигрывает, если в первом ходу забирает $1$ алмаз, а затем всегда действует так, чтобы сумма алмазов, взятых за ход Алекса и его собственный следующий ход, была равна $5$. Тогда при любом выборе Алекса Стив забирает последний алмаз.

\newpage

\section*{Задача 4}

У нас есть $13$ ворот, у каждого лампа «ВЫКЛ/ВКЛ». Изначально все лампы «ВЫКЛ». При подходе к воротам $i$ лампа сначала переключается, а потом:
\begin{itemize}
  \item если \emph{до переключения} была «ВЫКЛ», игрок телепортируется к первым воротам;
  \item если \emph{до переключения} была «ВКЛ», игрок проходит дальше.
\end{itemize}

Опишем состояние всех ламп как двоичное число из $13$ бит:
\[
b_1 b_2 \dots b_{13},\quad b_i = 0 \text{ для «ВЫКЛ», } b_i = 1 \text{ для «ВКЛ».}
\]
Изначально имеем $000\ldots0$ (все лампы «ВЫКЛ»).

Что делает одна попытка? Игрок идёт от ворот 1 вправо:
\begin{itemize}
  \item пока $b_i = 1$, он проходит через ворота $i$, а лампа переключается в $0$;
  \item на первом вороте с $b_j = 0$ он проваливается, лампа переключается в $1$, попытка заканчивается.
\end{itemize}

То есть каждый проход реализует ровно операцию «прибавить 1» к двоичному числу $b_1 b_2 \dots b_{13}$: хвост единиц превращается в нули, первый ноль превращается в единицу. Изначально число равно $0$, после первой попытки --- $1$, после второй --- $2$ и так далее.

Успешный проход возможен только тогда, когда \emph{перед попыткой} все лампы «ВКЛ», то есть состояние равно
\[
111\ldots1_2 = 2^{13}-1.
\]
В этой попытке игрок последовательно проходит все $13$ ворот и попадает в сокровищницу.

Значит, успешная попытка --- это $2^{13}$-я по счёту.

\subsection*{Ответ}
В бланк: $2^{13}$.
\newpage

\section*{Задача 5}

В туннель длиной $100$ блоков одновременно запустили $8$ вагонеток, которые не могут разъехаться. Каждая движется со скоростью $1$ блок в тик. Изначально:
\begin{itemize}
  \item Вагонетка 1: координата 10, вправо (R)
  \item Вагонетка 2: координата 25, вправо (R)
  \item Вагонетка 3: координата 30, влево (L)
  \item Вагонетка 4: координата 45, вправо (R)
  \item Вагонетка 5: координата 50, влево (L)
  \item Вагонетка 6: координата 70, влево (L)
  \item Вагонетка 7: координата 85, вправо (R)
  \item Вагонетка 8: координата 95, влево (L)
\end{itemize}
При лобовом столкновении две вагонетки мгновенно отскакивают и меняют направление, скорость остаётся той же. Как только вагонетка достигает координаты $0$ или $100$, она покидает туннель навсегда. Нужно найти:
\begin{enumerate}
  \item суммарное количество столкновений;
  \item через сколько тиков туннель покинет последняя вагонетка.
\end{enumerate}

Ключевое наблюдение: в задаче с упругими столкновениями \emph{одинаковых} тел, движущихся по прямой, можно мысленно заменить столкновение на «прозрачное прохождение» --- считать, что вагонетки просто проходят друг сквозь друга, не взаимодействуя. Тогда траектории остаются теми же, а времена ухода с концов и порядок покидания туннеля совпадают с реальной задачей. Но для подсчёта количества столкновений нам нужно чуть аккуратнее учесть встречу пар.

Сначала найдём время, через которое каждая вагонетка покинет туннель, если считать, что они не взаимодействуют (проходят друг сквозь друга):
\begin{itemize}
  \item Вагонетка 1: старт в 10, вправо, до 100: время $100 - 10 = 90$ тиков.
  \item Вагонетка 2: старт в 25, вправо, до 100: время $100 - 25 = 75$.
  \item Вагонетка 3: старт в 30, влево, до 0: время $30 - 0 = 30$.
  \item Вагонетка 4: старт в 45, вправо, до 100: время $100 - 45 = 55$.
  \item Вагонетка 5: старт в 50, влево, до 0: время $50 - 0 = 50$.
  \item Вагонетка 6: старт в 70, влево, до 0: время $70 - 0 = 70$.
  \item Вагонетка 7: старт в 85, вправо, до 100: время $100 - 85 = 15$.
  \item Вагонетка 8: старт в 95, влево, до 0: время $95 - 0 = 95$.
\end{itemize}

Максимальное время среди этих --- это $95$ тиков (у вагонетки 8). Значит, последней туннель покинет именно вагонетка 8, а время до полного опустошения туннеля равно $95$ тикам.

Теперь обсудим количество столкновений. Рассмотрим любую пару вагонеток. Если одна едет вправо, а другая --- влево, причём первая стоит левее второй (то есть их траектории на прямой пересекутся), то в реальной модели они обязательно \emph{один раз} столкнутся. В модели «прозрачного прохождения» их пути просто пересекаются в одной точке, но мы можем считать это столкновением.

Удобно считать столкновения по парам: для каждой пары вагонеток, стартующих с координат $x_i$ и $x_j$ ($x_i < x_j$), проверим:
\begin{itemize}
  \item если левая вагонетка едет вправо, а правая --- влево, то будет одно столкновение;
  \item в остальных случаях столкновения не будет.
\end{itemize}

Перечислим пары, где возможны столкновения.

\begin{itemize}
  \item Пара (2,3): вагонетка 2 (25,R), вагонетка 3 (30,L). Левая едет вправо, правая --- влево, будет 1 столкновение.
  \item Пара (4,5): 4 (45,R), 5 (50,L). Аналогично, 1 столкновение.
  \item Пара (4,6): 4 (45,R), 6 (70,L). Левая вправо, правая влево, траектории пересекутся, ещё 1 столкновение.
  \item Пара (7,8): 7 (85,R), 8 (95,L). Левая вправо, правая влево, 1 столкновение.
\end{itemize}

Проверим остальные пары для надёжности:
\begin{itemize}
  \item Пара (1,3): 1 (10,R), 3 (30,L) --- левая вправо, правая влево, тоже пересекутся, ещё 1 столкновение.
  \item Пара (1,5): 1 (10,R), 5 (50,L) --- аналогично, 1 столкновение.
  \item Пара (1,6): 1 (10,R), 6 (70,L) --- аналогично, 1 столкновение.
  \item Пара (2,5): 2 (25,R), 5 (50,L) --- аналогично, 1 столкновение.
  \item Пара (2,6): 2 (25,R), 6 (70,L) --- аналогично, 1 столкновение.
  \item Пара (3,4): 3 (30,L), 4 (45,R) --- левая едет \emph{влево}, правая вправо, они будут удаляться друг от друга, столкновения не будет.
  \item Пара (3,7): 3 (30,L), 7 (85,R) --- тоже расходятся.
  \item Пара (3,8): 3 (30,L), 8 (95,L) --- обе влево, не столкнутся.
  \item Пара (4,7): 4 (45,R), 7 (85,R) --- обе вправо, не столкнутся.
  \item Пара (5,6): 5 (50,L), 6 (70,L) --- обе влево.
  \item Пара (5,7): 5 (50,L), 7 (85,R) --- левая влево, правая вправо, расходятся.
  \item Пара (5,8): 5 (50,L), 8 (95,L) --- обе влево.
  \item Пара (6,7): 6 (70,L), 7 (85,R) --- левая влево, правая вправо, расходятся.
  \item Пара (6,8): 6 (70,L), 8 (95,L) --- обе влево.
\end{itemize}

Всего пар «левая вправо, правая влево»:
\[
(1,3), (1,5), (1,6), (2,3), (2,5), (2,6), (4,5), (4,6), (7,8).
\]
Это $9$ пар, значит происходит $9$ столкновений.

Итого:
\begin{itemize}
  \item суммарное количество столкновений --- $9$;
  \item последняя вагонетка покинет туннель через $95$ тиков.
\end{itemize}

\subsection*{Ответ}
В бланк: $9\ 95$.\\
Пояснение: время ухода вагонеток равно времени, за которое каждая бы дошла до конца без столкновений; максимальное из этих времён --- $95$ тиков. Столкновения считаем по парам вагонеток: каждая пара «левая вправо, правая влево» даёт ровно одно столкновение; таких пар $9$.

\end{document}